\subsection{Continuous-time Markov chains}
\label{m:sec:ctmc}

\subsubsection{Exponential clocks}
\label{m:sec:exp}

Every continuous-time model in this thesis is built from events whose rate of
occurrence does not depend on how long the system has been waiting for them.
That single assumption forces the exponential distribution: a \rv{} $T$ is
exponentially distributed with rate $\theta>0$, written
$T\sim\mathrm{Exp}(\theta)$, if it has density $\theta\ee^{-\theta t}$ on
$t\ge0$, mean $1/\theta$ and variance $1/\theta^2$. The property that matters is
memorylessness: for $t_1,t_2>0$,
\begin{equation}
\label{m:eq:memoryless}
\pr{T > t_1+t_2 \mid T > t_1}
= \frac{\ee^{-\theta(t_1+t_2)}}{\ee^{-\theta t_1}}
= \ee^{-\theta t_2}
= \pr{T > t_2},
\end{equation}
so a clock that has already run for time $t_1$ is statistically
indistinguishable from a fresh one. The exponential law is the only continuous
distribution with this property, which is why a process assembled from
constant-rate events depends on its present state alone and not on its history.

When several events compete, the following result organises everything that
follows. Let $T_1,\dots,T_m$ be independent with $T_i\sim\mathrm{Exp}(\theta_i)$,
and let $T=\min_iT_i$ be the time of the first event.

\begin{proposition}[Competing clocks]
\label{m:prop:race}
Writing $\Theta=\sum_i\theta_i$,
\begin{equation}
\label{m:eq:race}
T \sim \mathrm{Exp}(\Theta),
\qquad
\pr{T_k = T} = \frac{\theta_k}{\Theta},
\end{equation}
and the identity of the winning clock is independent of $T$.
\end{proposition}

\begin{proof}
For the first claim, $\pr{T>t}=\prod_i\pr{T_i>t}=\ee^{-\Theta t}$. For the
second, condition on the value of $T_k$:
\begin{equation*}
\pr{T_k = T}
= \int_0^\infty \theta_k\ee^{-\theta_k u}\prod_{i\neq k}\ee^{-\theta_i u}\,\dd u
= \theta_k\int_0^\infty \ee^{-\Theta u}\,\dd u
= \frac{\theta_k}{\Theta},
\end{equation*}
and truncating the integral at $t$ gives
$\pr{T_k=T,\ T>t}=(\theta_k/\Theta)\ee^{-\Theta t}
=\pr{T_k=T}\pr{T>t}$, which is the asserted independence.
\end{proof}

\Cref{m:prop:race} converts rates into probabilities as soon as one asks
\emph{which} of several things happened rather than when, and it is the whole
content of the direct simulation method: draw the next event time from
$\mathrm{Exp}(\Theta)$, draw the event $k$ with probability $\theta_k/\Theta$,
advance and repeat \cite{Gillespie1977}. Because the two draws are independent
and the exponential law is memoryless, the resulting trajectory has exactly the
law of the process. The scheme returns in \cref{m:sec:coupled}, where the rates
themselves move between jumps.

\subsubsection{Generators and forward equations}
\label{m:sec:generator}

Let $\{X_t:t\ge0\}$ take values in a countable state space $\mathcal S$, which
throughout will be a set of population counts or of tuples of counts. The process
is a \emph{continuous-time Markov chain} if, for all $s,t\ge0$,
\begin{equation}
\label{m:eq:markov}
\pr{X_{t+s} = j \mid X_t = i,\ \{X_u\}_{u<t}} = \pr{X_{t+s}=j\mid X_t = i},
\end{equation}
and \emph{time homogeneous} when the right-hand side depends on $s$ but not on
$t$; time-inhomogeneous chains are the subject of \cref{m:sec:inhomog}. A
homogeneous chain is specified by its \emph{generator} $Q=(q_{ij})$, whose
off-diagonal entries are transition rates,
\begin{equation}
\label{m:eq:generator}
q_{ij} = \lim_{\Delta t\to0}\frac{p_{ij}(\Delta t)}{\Delta t}\quad(j\neq i),
\qquad
q_{ii} = -\sum_{j\neq i}q_{ij} =: -q_i ,
\end{equation}
so that every row sums to zero. Equivalently,
\begin{equation}
\label{m:eq:infinitesimal}
\pr{X_{t+\Delta t} = j \mid X_t = i} = q_{ij}\,\Delta t + o(\Delta t)
\qquad (j \neq i),
\end{equation}
with complementary probability $1-q_i\Delta t+o(\Delta t)$ of no transition. By
\cref{m:prop:race} the chain leaves state $i$ after a holding time distributed
as $\mathrm{Exp}(q_i)$ and then jumps to $j$ with probability $q_{ij}/q_i$. A
state with $q_i=0$ is \emph{absorbing}: the chain that reaches it stays there
forever. Extinction and rupture are both modelled as absorbing states, and the
conditional methods of \cref{m:sec:methct} concern the behaviour of a chain
conditioned on not yet having been absorbed.

Differentiating the Chapman--Kolmogorov identity gives the \emph{forward
Kolmogorov equations},
\begin{equation}
\label{m:eq:kolmogorov}
\ddt{p_{ij}(t)} = \sum_{k} p_{ik}(t)\,q_{kj},
\qquad\text{that is}\qquad
\dda{P}{t} = P\,Q ,
\end{equation}
which for a fixed initial state and with $p_j(t)=\pr{X_t=j}$ become the master
equation
\begin{equation}
\label{m:eq:master}
\ddt{p_j(t)} = \sum_{k\neq j}q_{kj}\,p_k(t) - q_j\,p_j(t):
\end{equation}
probability flows into state $j$ from every state that can reach it and out of
$j$ at total rate $q_j$. Every model in this thesis is specified by writing down
\cref{m:eq:infinitesimal} and reading off \cref{m:eq:master}; the companion
\emph{backward} equations and their uses are developed in
\cref{m:sec:backward}.

\subsubsection{The Poisson process}
\label{m:sec:poisson}

The Poisson process is the CTMC from which all others are built, and the
prototype of a pure-birth process. Its state space is $\{0,1,2,\dots\}$ and its
only transitions are $n\to n+1$ at constant rate $\lambda$, so
$q_{n,n+1}=\lambda$ and $q_n=\lambda$ for every $n$: the holding time in every
state is $\mathrm{Exp}(\lambda)$, independent of everything else by
\cref{m:eq:memoryless}, and the process counts the arrivals of a stream of
events occurring one at a time at constant rate.

The master equation \cref{m:eq:master} reads
\begin{equation*}
\ddt{p_n(t)} = \lambda p_{n-1}(t) - \lambda p_n(t),
\qquad
p_n(0) = \ind_{\{n=0\}},
\end{equation*}
which solves inductively to the Poisson law
\begin{equation}
\label{m:eq:poissonlaw}
N_t \sim \mathrm{Poisson}(\lambda t),
\qquad
\pr{N_t = n} = \ee^{-\lambda t}\frac{(\lambda t)^n}{n!},
\end{equation}
with $\ex{N_t}=\var{N_t}=\lambda t$. Two immediate readings of
\cref{m:eq:poissonlaw} are worth keeping. The number of events in disjoint time
intervals are independent \rv s, because the exponential holding times restart
afresh after every event; and the jump times of any time-homogeneous CTMC whose
exit rate happens to stay constant --- for instance a pure-birth process before
its first death --- form a Poisson process.

The Poisson process is also the limit object behind thinning and superposition:
retaining each event of a rate-$\lambda$ process independently with probability
$p$ leaves a Poisson process of rate $p\lambda$, and independent Poisson streams
of rates $\lambda_1,\lambda_2$ combine to one of rate $\lambda_1+\lambda_2$ ---
both direct consequences of \cref{m:eq:poissonlaw} and independence. These facts
are used quietly below whenever a stream of events is split by type or merged.

\subsubsection{Birth--death processes}
\label{m:sec:linbd}

The linear birth--death process is the continuous-time counterpart of the
Galton--Watson process and the running example of \cref{m:sec:methct}. Let
$X_t$ count individuals, each of which independently gives birth at rate
$\lambda$ and dies at rate $\mu$. The non-zero rates out of state $n$ are
\begin{equation}
\label{m:eq:bdrates}
q_{n,n+1} = \lambda n,
\qquad
q_{n,n-1} = \mu n ,
\end{equation}
so that $q_n=(\lambda+\mu)n$ and $n=0$ is absorbing: the rates are linear in the
count, which is what makes the process exactly tractable by the methods of
\cref{m:sec:methct}. The master equation reads
\begin{equation}
\label{m:eq:bdmaster}
\ddt{p_n(t)} = \lambda(n-1)p_{n-1}(t) + \mu(n+1)p_{n+1}(t) - (\lambda+\mu)n\,p_n(t).
\end{equation}
With the \pgf{} $G(z,t)=\ex{z^{X_t}}$, multiplying \cref{m:eq:bdmaster} by
$z^n$, summing over $n$ and shifting indices collapses the infinite family of
state probabilities to a single partial differential equation,
\begin{equation}
\label{m:eq:bdpde}
\pddt{G} = \bigl(\lambda z^2 - (\lambda+\mu)z + \mu\bigr)\pdd{G}{z}
= (\lambda z-\mu)(z-1)\pdd{G}{z},
\qquad G(z,0) = z^{N},
\end{equation}
for a process started from $N$ individuals. This conversion --- master equation
to first-order generating-function equation --- is the template for every
generating-function calculation in this thesis, and its solution by the method of
characteristics is prepared in \cref{m:sec:moc}. The result, obtained there, is
Kendall's formula \cite{Kendall1948}:
\begin{equation}
\label{m:eq:bdsolution}
G(z,t) = \lrb{\frac{\mu(1-z) - (\mu-\lambda z)\ee^{-(\lambda-\mu)t}}
{\lambda(1-z) - (\mu-\lambda z)\ee^{-(\lambda-\mu)t}}}^{\!N},
\qquad \lambda\neq\mu ,
\end{equation}
from which the mean $\ex{X_t}=N\ee^{(\lambda-\mu)t}$ and the extinction
probability by time $t$,
\begin{equation}
\label{m:eq:p0t}
p_0(t) =
\begin{cases}
\displaystyle\lrb{\frac{\mu-\mu\ee^{(\mu-\lambda)t}}{\lambda-\mu\ee^{(\mu-\lambda)t}}}^{\!N},
& \lambda\neq\mu,\\[14pt]
\displaystyle\lrb{\frac{\lambda t}{1+\lambda t}}^{\!N},
& \lambda=\mu,
\end{cases}
\end{equation}
follow by differentiation and by setting $z=0$. Letting $t\to\infty$ recovers
the familiar threshold: extinction is certain when $\lambda\le\mu$ and occurs
with probability $(\mu/\lambda)^N$ when $\lambda>\mu$. The quadratic coefficient
in \cref{m:eq:bdpde} factorises with roots $z=1$ and $z=\mu/\lambda$; the same
factorisation, in a different variable, reappears in \cref{m:app:moc}.

A useful extension, which later chapters use, lets the rates vary slowly with the
population itself. If $\lambda,\mu$ are replaced by $\lambda(n),\mu(n)$ --- for
instance $\lambda(n)=\lambda_0(1-n/K)$ to model crowding --- the chain remains
Markov but is no longer exactly solvable by the closed route above; mean-field
treatments replace $n$ by $\ex{X_t}$ and recover deterministic density-dependent
equations \cite{Allen2010}. The time-inhomogeneous linear case, in which the
rates are prescribed functions of $t$ alone, remains tractable and is taken up
next.

\subsubsection{Birth--death--catastrophe processes}
\label{m:sec:ctbdc}

Alongside extinction by death, several later chapters need a second absorbing
mechanism: an abrupt event that removes the entire population at once. The
continuous-time birth--death--catastrophe process adds that channel to
\cref{m:sec:linbd}.

\begin{definition}[Birth--death--catastrophe process]
\label{m:def:bdc}
Let $X_t$ count individuals in a compartment. Each individual independently gives
birth at rate $\lambda$ and dies at rate $\mu$, and the compartment additionally
undergoes a catastrophe --- a rupture that removes the entire population at once
--- at a rate $\rho$ per individual, so that the total catastrophe rate in state
$n$ is $\rho n$. Writing $\dagger$ for the post-catastrophe state, the non-zero
transition rates out of $n\ge1$ are
\begin{equation}
\label{m:eq:bdcrates}
q_{n,n+1} = \lambda n,
\qquad
q_{n,n-1} = \mu n,
\qquad
q_{n,\dagger} = \rho n ,
\end{equation}
and both $n=0$ and $\dagger$ are absorbing.
\end{definition}

Two features distinguish \cref{m:def:bdc} from \cref{m:sec:linbd}, and both
matter later. The process has \emph{two} absorbing states, reached by mechanisms
that are distinguishable and not interchangeable: extinction empties the
compartment gradually, catastrophe empties it in one event and releases its
contents. And the catastrophe rate is proportional to the population, so the
compartment becomes more likely to rupture the fuller it is.

Restricting attention to realisations that have not yet ruptured and writing
$G(z,t) = \ex{z^{X_t}\ind\{X_t\neq\dagger\}}$, the master equation for
\cref{m:eq:bdcrates} yields
\begin{equation}
\label{m:eq:bdcpde}
\pddt{G} = \bigl(\lambda z^2-(\lambda+\mu+\rho)z+\mu\bigr)\pdd{G}{z},
\qquad G(z,0) = z^{N},
\end{equation}
which differs from \cref{m:eq:bdpde} only in the coefficient of $z$. The
quadratic no longer factorises with a root at $z=1$, and that single change is
responsible for most of what is interesting about the process: probability is
not conserved on the non-ruptured states, and $G(1,t)$ is the survival
probability rather than one. The discrete-generation counterpart, its exact
no-rupture identities, and the killed-generator formulation of
quasi-stationarity for this process are developed in \cref{m:app:dbdc}; rates,
moments and state probabilities for \cref{m:def:bdc} itself are treated
\fwd{bdc}{in the later chapters on birth--death--catastrophe processes}, with
the two-type extension \fwd{multitype}{in the chapter on multi-type processes}
and the release of a ruptured compartment's contents into a shared medium
\fwd{rupture}{in the chapter on compartment rupture}.

\FloatBarrier
